% ==========================================================
% HyperSLR Methods Snippet — auto-generated by Cell 41
% Experiment: hypersrl_20260312_112009
% Provenance hash: efe9b1182ae9def1c1e5227bf76f2b58
% Generated: 20260312_112009
% ==========================================================

\subsection{Data and Study Region}
The HyperSLR framework ingests tide gauge records from the Permanent
Service for Mean Sea Level (PSMSL) for three index stations representing
Bangladesh (stations 369, 543, 1270), Vietnam (841, 1449), and the
Maldives (1707). The merged analysis dataset comprises 2,644 monthly
observations spanning 1950-01 to 2024-12, standardised to a common
2022-12-31 anomaly baseline. Eight predictive features are retained from
the Cell 28 PINN architecture: sea-level anomaly (slrv2), tide-gauge
residual, ENSO index, AVISO satellite altimetry, vertical land motion
(VLM), CMIP6 thermosteric expansion, CMIP6 glacier contribution, and
sea-level momentum. Four engineered features are added: monthly
harmonic encodings ($\sin$/$\cos$) and an ENSO$\times$AVISO interaction
term.

\subsection{Physics-Informed Neural Network Architecture}
The base PINN (Cell 28) follows a dual-head LSTM architecture with
input size 8, hidden size 128, two recurrent layers with dropout
$p = 0.2$, a primary regression head
(Linear(128, 64) $\to$ ReLU $\to$ Dropout(0.2) $\to$ Linear(64, 1)),
and a physics residual head (Linear(128, 32) $\to$ ReLU $\to$
Linear(32, 1)). Total trainable parameters: 215,234.
Architecture content hash (MD5): \texttt{bb8441ca}.

\subsection{Continual Learning Protocol}
Model adaptation follows an elastic-weight-consolidation-inspired
continual learning protocol. The LSTM encoder is frozen
(\texttt{lstm})
and only the regression heads are fine-tuned on new observations.
Adaptation is triggered when a minimum of None
new monthly observations become available beyond the baseline date.
Training uses the Adam optimiser with learning rate $\eta = 0.0001$,
batch size 16, sequence length 12 months, physics loss
weight $\lambda_{\mathrm{phys}} = 0.4$, maximum 30
epochs, and early stopping with patience 10. A quality gate
rejects updates where the post-adaptation RMSE exceeds the
pre-adaptation RMSE by more than 5\%.

\subsection{Risk Threshold Detection}
Cumulative sea-level rise is computed relative to the 2022-12-31
baseline using a $\pm$6-month windowed mean. Risk thresholds are set
at physically motivated levels reflecting national coastal
vulnerability indices (CVI): Bangladesh
(AMBER: 150\,mm, RED: 250\,mm, CRITICAL: 350\,mm;
CVI\,=\,0.847, VLM\,=\,$-2.0$\,mm\,yr$^{-1}$),
Vietnam
(AMBER: 100\,mm, RED: 200\,mm, CRITICAL: 300\,mm;
CVI\,=\,0.793, VLM\,=\,$-8.9$\,mm\,yr$^{-1}$), and
Maldives
(AMBER: 50\,mm, RED: 80\,mm, CRITICAL: 100\,mm;
CVI\,=\,0.921, VLM\,=\,0.0\,mm\,yr$^{-1}$).
Compound flood probabilities are estimated via a Gumbel copula
($\theta = 2.0$) coupling SSP2-4.5 SLR exceedance with
historically-fitted storm surge return periods.
At the time of this analysis, the system risk level is
\textbf{GREEN} with Bangladesh recording
108.677\,mm cumulative rise above baseline.

\subsection{Software and Reproducibility}
All analyses were conducted in Python 3.12.12 on CPU only
using PyTorch 2.10.0+cpu within a Google Colab
environment. Complete pipeline reproducibility is ensured via
content-addressed MD5 hashing of all cell source files
(provenance hash: \texttt{efe9b1182ae9def1}\ldots).
The full pipeline is orchestrated by Cell 40 with
APScheduler-based scheduling and halt-on-failure semantics.
Code and data are version-controlled via the HyperSLR
experiment registry (Cell 41).
